\documentclass[../main.tex]{subfiles}

\begin{document}
    \subsection{Bagi THR}
    Holil bingung bagaimana cara membagikan THR-nya. Sehingga ia berencana membuat sebuah program yang dapat menangani masalah tersebut. Pertama, Holil ingin mendata setiap mahasiswa yang akan ia berikan THR dan mengelompokkan mereka berdasarkan tipe mahasiswa seperti apa mereka, Holil lalu membagi tipe mahasiswa menjadi 3 tipe yaitu mahsiswa kupu-kupu, kura-kura dan kuda-kuda.

    \inputfigure[n]{res/m2i1.png}{0.3}{Keluaran Pertama Program Bagi THR}{fig:m2:i1}
    Selanjutnya, Nominal THR yang Holil berikan berbeda setiap tipe mahasiswa, di antaranya mahasiswa kupu-kupu menerima Rp.500K/org, kura-kura Rp.850K/org dan kuda-kuda 1jt/org. Sebelum memberikan THR, Holil menanyakan durasi Jam kuliah mahasiswa selama Ramadhan yang nantinya Holil konversi ke satuan detik dan akan ia tuliskan di amplop pemeberian THR kepada mahasiswa. Setelah beberapa saat, Holil berpikir ingin memberikan tambahan THR untuk mahasiswa yang memiliki kuliah malam, sehingga Holil juga menanyakan durasi kuliah malam dalam Detik.

    \inputfigure[n]{res/m2i2.png}{0.35}{Keluaran Kedua Program Bagi THR}{fig:m2:i2}
    Mahasiswa dengan jam kuliah malam antara 1 – 2 jam ditambah Rp.200K, antara 2 – 3 jam ditambah Rp.500K, > 3 jam akan dikali 2 jumlah THR awal. Lalu, Holil menyiapkan amplop baru untuk mahasiswa yang memiliki kuliah malam, Holil juga menuliskan total durasi kuliah berdasarkan jam, menit, detik di amplop baru tersebut.

    \inputfigure[n]{res/m2i3.png}{0.1}{Keluaran Ketiga Program Bagi THR}{fig:m2:i3}
    Namun, beberapa hari ini Holil sibuk karena harus mencari baju untuk lebaran dan Holil tak sempat menyelesaikan program THR tersebut. Maka dari itu, agar kalian segera mendapatkan THR dari Holil, bantulah Holil membuat program THR sesuai keinginannya.

    \subsection{Motif Baju}
    Holil meminta dibuatkan 3 jenis motif baju berbeda yakni motif bagun datar, pohon cemara dan motif pola bilangan.
    \inputfigure[n]{res/m2i4.png}{0.16}{Keluaran Pertama Program Motif Baju}{fig:m2:i4}

    Dalam motif bagun datar, Holil memiliki 3 keinginan yakni motif dengan bentuk persegi, segitiga dan persegi lagi. Motif Persegi terbuat dari pola sudut siku huruf A – Z yang semakin mengecil dari yang kiri atas ke kanan bawah sesuai ukuran yang di-{\ioi}-kan. Segitiga terdiri dari kombinasi huruf dan angka yang saling bergantian. Apabila {\ioi} yang dimasukkan ganjil, maka akan dimulai dari huruf B sedangkan apabila {\ioi} genap akan dimulai dari A. Lalu Motif Persegi Lagi yang terdiri dari pola bilangan bersambung dari kiri ke kanan dan sebaliknya pada baris berikutnya.
    \inputfigure[n]{res/m2i5.png}{0.16}{Keluaran Kedua Program Motif Baju}{fig:m2:i5}

    Motif selanjutnya adalah Motif Pohon Cemara yang terdiri dari kumpulan karakter ``*''. Bentuk dari pohon ini merupakan kombinasi dari pola segitiga dengan ukuran tergantung tinggi dan lebar yang diminta, yang bertumpuk ke bawah dengan pola tumpukan +2 dari ukuran lebar awal segitiga teratas dan pola kotak yang memiliki ukuran yang tetap sebagai batangnya pada bagian bawah.
    \inputfigure[n]{res/m2i6.png}{0.16}{Keluaran Ketiga Program Motif Baju}{fig:m2:i6}

    Motif berikutnya berupa pola bilangan yang outputnya dipengaruhi oleh {\ioi}. Apabila {\ioi}-an bernilai ganjil, maka program akan menampilkan pola bilangan seperti gambar di samping kiri. Apabila {\ioi} bernilai genap, maka program akan menampilkan pola bilangan seperti gambar di samping kanan.
    \inputfigure[n]{res/m2i7.png}{0.3}{Keluaran Keempat Program Motif Baju}{fig:m2:i7}

    Menu Program Pola bilangan pada {\ioi}-an genap akan menampilkan pola bilangan prima, di mana jika {\ioi} yang dimasukkan kurang dari sama dengan 50, maka program akan mendeteksi bilangan yang mengandung angka ``3'' kemudian bilangan tersebut akan dikalikan dengan n ({\ioi}). Apabila {\ioi} melebihi 50, maka akan ditampilkan pola bilangan prima biasa.
    \inputfigure[n]{res/m2i8.png}{0.3}{Keluaran Kelima Program Motif Baju}{fig:m2:i8}

    \subsection{Hilal}
    Holil berinisiatif untuk membuat program yang bisa digunakan untuk memantau Hilal, untuk menentukan awal bulan Ramadan dan Idul Fitri. Dengan program ini, Holil dapat mengidentifikasi fase bulan, yang akan membantunya merencanakan produksi dengan lebih tepat. Dengan demikian, ia bisa memastikan bahwa produksinya selesai tepat waktu untuk di- launching saat hari raya, sesuai dengan targetnya. Bantulah Holil membuat program yang dapat menampilkan pergerakan fase bulan, dengan menggunakan persamaan umum lingkaran x$^2$ + y$^2$ = r$^2$, untuk menampilkan bulannya.
    \inputfigure[n]{res/m2i9.png}{1.0}{Gambar Bulan Yang Akan Dianimasikan Oleh Program Hilal}{fig:m2:i9}
\end{document}
